% -*- coding=utf-8 -*-
\chapter{预备知识}  % \chapter{短标题}{长标题} 页眉会显示短标题，正文会显示长标题，“\\"会强制换行，”\ “会显示空格
我是预备知识

在离散数学中，图（Graph）是用于表示物体与物体之间存在某种关系的结构。数学抽象后的“物体”称作节点或顶点（英语：Vertex，node或point），节点间的相关关系则称作边。[1]在描绘一张图的时候，通常用一组点或小圆圈表示节点，其间的边则使用直线或曲线。

用$\mathbb{R}$和$\mathbb{R}^{n}$分别表示实数域和$n$维实数域. 对任意的实数$x \in \mathbb{R}$, $\lfloor x \rfloor$是向下取整函数，表示不大于$x$的最大整数，且$ \lfloor x \rfloor \leq x$；$\lceil x \rceil$是向上取整函数，表示不小于$x$的最小整数，且$\lceil x \rceil \geq x$.
 
\section{向量、矩阵和范数}
% https://oi-wiki.org/math/linear-algebra/vector-space/#%E7%9B%B4%E8%A7%82%E7%90%86%E8%A7%A3
二维平面上的一点$P=(x,y)^{\top}$

\begin{lemma}\label{Lemma4EmYqK}
	假设$e_p,~p=1,2,\cdots,m$是$m$维单位矩阵的第$p$列和列向量$Y_{q} = \left( y_{1,q}, y_{2,q}, \cdots, y_{m,q} \right)^\top \in \mathbb{R}^{m}$, 且$y_{p,q} \neq 0$,  则
	\begin{equation*}
	 \left( e_1, \cdots,e_{p-1}, Y_{q}, e_{p+1}, \cdots, e_m \right)^{-1}  = 
	  	   		\begin{pmatrix} 
		  		1 	   & 0      & \cdots  & 0 			  & -\frac{y_{1,q}}{y_{p,q}}  	& 0 	& \cdots   & 0\\
		  		0 	   & 1      & \cdots  & 0 			  & -\frac{y_{2,q}}{y_{p,q}}  	& 0 	& \cdots   & 0\\
		  		\vdots & \vdots & \ddots  & \vdots 	      &  \vdots     	& \vdots  	& \ddots  & \vdots\\
		  		0      & 0      & \cdots  & 1  			  & -\frac{y_{p-1,q}}{y_{p,q}}  	& 0 	& \cdots   & 0\\	  		
		  		0      & 0      & \cdots  & 0  			  & \frac{1}{y_{p,q}}  	& 0 	& \cdots   & 0\\
		  		0      & 0      & \cdots  & 0  			  & -\frac{y_{p+1,q}}{y_{p,q}}  	& 1 	& \cdots   & 0\\
		  		\vdots & \vdots & \ddots  & \vdots  	  & \vdots  		& \vdots  	& \ddots   & \vdots\\	
		  		0      & 0      & \cdots  & 0  			   & -\frac{y_{m,q}}{y_{p,q}}     & 0     & \cdots   & 1\\	  		
			\end{pmatrix} .
	\end{equation*}
\end{lemma}



 \section{开集和闭集}
以$x \in \mathbb{R}^{n}$为球心，$r > 0$为半径的开球记为 
 \begin{equation*}
     \mathcal{B}\left(x,r\right) \triangleq \left\{ y \in \mathbb{R}^{n} : \Vert y - x \Vert_{2} < r \right\}
 \end{equation*}
 和相应的闭球
  \begin{equation*}
     \overline{\mathcal{B}}\left(x,r\right) \triangleq \left\{ y \in \mathbb{R}^{n} : \Vert y - x \Vert_{2} \leq r \right\}.
 \end{equation*}
\begin{definition}\label{Chapter0Open}
    设非空集合$\mathcal{S} \subseteq \mathbb{R}^{n}$. 若对任意的$x \in \mathcal{S}$, 总存在$r > 0$, 使得$\mathcal{B}\left(x,r\right) \subseteq \mathcal{S}$, 则称$\mathcal{S}$为\textbf{开集}. \textbf{补集}$~\mathcal{S}^{\mathrm{c}} \triangleq \left\{ x \in \mathbb{R}^{n} : x \notin \mathcal{S} \right\}$为开集的集合$\mathcal{S}$称为\textbf{闭集}，其中\textbf{补集}$~\mathcal{S}^{\mathrm{c}}$也被记为$~\overline{\mathcal{S}}$.
\end{definition}
 显然地，开球$\mathcal{B}\left(x,r\right)$和闭球$\overline{\mathcal{B}}\left(x,r\right)$分别是开集和闭集. 包含点$x \in \mathbb{R}^{n}$的开集称为$x$的一个\textbf{领域}. 全空间$\mathbb{R}^{n}$和空集$\emptyset$即是开集又是闭集. 进一步地，
 \begin{itemize}
  \setlength{\itemsep}{0pt}
\setlength{\parsep}{0pt}
\setlength{\parskip}{0pt}
     \item[$(1)$] 有限个开集$S_i(i=1,\cdots,m)$的交集$\bigcap_{i=1}^{m}S_i$是开集，任意个开集$S_i(i \in \mathcal{I} )$的并集$\bigcup_{i \in \mathcal{I}}S_i$是开集，其中$\mathcal{I}$为有限或无限的指标集合；
     \item[$(2)$] 有限个闭集$S_i(i=1,\cdots,m)$的并集$\bigcup_{i=1}^{m}S_i$是闭集，任意个闭集$S_i(i \in \mathcal{I} )$的交集$\bigcap_{i \in \mathcal{I}}S_i$是闭集，其中$\mathcal{I}$为有限或无限的指标集合.
 \end{itemize}

给定一个非空集合$\mathcal{S} \subseteq \mathbb{R}^{n}$与点$x \in \mathcal{S}$, 若存在$r > 0$, 使得$\mathcal{B}\left(x,r\right) \subseteq \mathcal{S}$, 则称$x$为$\mathcal{S}$的\textbf{内点}. 非空集合$\mathcal{S}$中所有\textbf{内点}构成的集合称为$\mathcal{S}$的\textbf{内部},记为$\rm{int} \mathcal{S}$. 显然，$\rm{int} \mathcal{S}$是开集，这是因为在$\rm{int} \mathcal{S}$中的任何一个点，都可找到一个属于$\rm{int} \mathcal{S}$的开球. 另外，称包含$\mathcal{S}$的最小闭集，即所有包含$\mathcal{S}$的闭集的交集，为$\mathcal{S}$的\textbf{闭包}, 记为$\rm{cl} \mathcal{S}$. 属于$\rm{cl} \mathcal{S}$但不属于$\rm{int} \mathcal{S}$的所有点构成的集合称为$\rm{int} \mathcal{S}$的边界，记为$\rm{bd} \mathcal{S}$. 如果存在一个充分大的数$r > 0$, 使得$\mathcal{S} \subseteq \mathcal{B}\left(0,r\right)$, 则称$\mathcal{S}$是\textbf{有界}. 进一步地，称有\textbf{有界闭集}$\mathcal{S}$为\textbf{紧集}. 注意，此处的\textbf{紧集}的定义只适合有限维空间.



\section{凸集和凸函数}
图中的边可以是有方向或没有方向的。例如在一张图中，如果节点表示聚会上的人，而边表示两人曾经握手，则该图就是没有方向的，因为甲和乙握过手也意味着乙一定和甲握过手。相反，如果一条从甲到乙的边表示甲欠乙的钱，则该图就是有方向的，因为“曾经欠钱”这个关系不一定是双向的。前一种图称为无向图，后一种称为有向图。

\begin{definition}\label{Chapter0Cone}
    设非空集合$\mathcal{C} \subseteq \mathbb{R}^{n}$. 如果满足对任意的$x \in \mathcal{C}$和$\alpha \in [0,+\infty)$, 都有$\alpha x \in \mathcal{C}$, 则称$\mathcal{C}$为\textbf{锥}.
\end{definition}

从锥的定义 \ref{Chapter0Cone} 可知，锥$\mathcal{C}$是包含所有以原点为起点，并通过$\mathcal{C}$中一点的射线的集合. 当锥$\mathcal{C}$是凸集，称之为\textbf{凸锥}；当锥$\mathcal{C}$是闭集时，称之为\textbf{闭锥}；当锥$\mathcal{C}$即是凸集又是闭集时，则称之为\textbf{闭凸锥}.

\section{梯度、Hessian矩阵和Jacobi矩阵}

\begin{definition}[梯度\index{梯度}]
    给定函数 $f: \mathbb{R}^n \to \mathbb{R}$，且$f$在点$x$的一个邻域内有定义，若存在向量 $g \in \mathbb{R}^n$ 满足
	\begin{equation*}
	   \lim_{p \to 0} \frac{f(x+p) - f(x) - g^{\mathrm{T}} p}{\|p\|} = 0, 
	\end{equation*}
    其中 $\| \cdot \|$ 是任意的向量范数，就称 $f$ 在点 $x$ 处\textup{Fr\'{e}chet} 可微．此时 $g$ 称为 $f$ 在点 $x$ 处的\textbf{梯度}，记作 $\nabla f(x)$．
		如果对区域 $D$ 上的每一个点 $x$ 都有 $\nabla f(x)$ 存在，则称 $f$ 在 $D$ 上\textup{Fr\'{e}chet}可微．\index{Fr\'{e}chet 可微}
\end{definition}

若 $f$ 在点 $x$ 处的梯度存在，在定义式中令 $p = \varepsilon e_i$，$e_i$是第$i$个分量为$1$的单位向量，可知 $\nabla f(x)$ 的第 $i$ 个分量为 $\dfrac{\partial f(x)}{\partial x_i}$．因此 
	\begin{equation*}
	   \nabla f(x) = \left[ \frac{\partial f(x)}{\partial x_1}, \frac{\partial f(x)}{\partial x_2}, \cdots, \frac{\partial f(x)}{\partial x_n} \right]^{\mathrm{T}}.
	\end{equation*}
 
\begin{definition}[Hessian矩阵]
	如果函数$f(x): \mathbb{R}^n \to \mathbb{R}$在点$x$处的二阶偏导数
		$\frac{\partial^2 f(x)}{\partial x_i \partial x_j}$ $i,j
		=1,2,\cdots,n$都存在，则$f$ 在点 $x$ 处的\rm{Hessian}矩阵为:
		\[ 
		\mathbf{H} = \mathbf{H}_f \left( x\right)  \triangleq \nabla^2 f(x) =
		\begin{bmatrix}
			\dfrac{\partial^2 f(x)}{\partial x_1^2}            & \dfrac{\partial^2 f(x)}{\partial x_1 \partial x_2} & \dfrac{\partial^2 f(x)}{\partial x_1 \partial x_3} & \cdots & \dfrac{\partial^2 f(x)}{\partial x_1 \partial x_n} \\
			\dfrac{\partial^2 f(x)}{\partial x_2 \partial x_1} & \dfrac{\partial^2 f(x)}{\partial x_2^2}            & \dfrac{\partial^2 f(x)}{\partial x_2 \partial x_3} & \cdots & \dfrac{\partial^2 f(x)}{\partial x_2 \partial x_n} \\
			\vdots                                             & \vdots                                             & \vdots                                             &        & \vdots                                             \\
			\dfrac{\partial^2 f(x)}{\partial x_n \partial x_1} & \dfrac{\partial^2 f(x)}{\partial x_n \partial x_2} & \dfrac{\partial^2 f(x)}{\partial x_n \partial x_3} & \cdots & \dfrac{\partial^2 f(x)}{\partial x_n^2}            
		\end{bmatrix}
		 \in \mathbb{R}^{n \times n}.
		\]
		特别地，Hessian矩阵的第$i$行第$j$列的元素为： $ \mathbf{H}_{ij} = \left(\mathbf{H}_f\left( x\right)  \right)_{ij} = \dfrac{\partial^2 f(x)}{\partial x_i \partial x_j}.$
\end{definition}